\section{Every height at large widths}\label{sec:general}

The constructions of Sections~\ref{sec:four} and~\ref{sec:five} were found one height at a time.
For large widths a single pattern attains the row bound at every height. Call a singleton column
a \emph{sea} when its black point lies in row $1$ or in row $h$.

\begin{theorem}[large widths]\label{thm:general}
Let $h\ge3$ and put $M_h=q_h+(q_h-1)(2W_h+1)$. Then $f(h,m)=m+q_h$ for every $m\ge M_h$.
\end{theorem}

The upper bound is Theorem~\ref{thm:rowbound}. The lower bound follows from the next proposition,
which needs only a triangle-free graph on the rows.

\begin{proposition}[alternating seas]\label{prop:seas}
Let $h\ge3$ and let $H$ be a triangle-free graph on the rows with edges $e_1,\dots,e_q$.
Let $Z$ be the blackout of the $h\times m$ grid with doubles in columns $p_1<\dots<p_q$, where
column $p_j$ is black exactly in the two rows of $e_j$ and $p_{j+1}-p_j\ge2W_h+2$, and with a sea
in every other column. The seas to the left of $p_1$ are black in row $1$, and the row of the
seas switches between $1$ and $h$ at every double. Then $Z$ is valid and $|Z|=m+q$.
\end{proposition}

\begin{proof}[Proof of Theorem~\ref{thm:general}]
Take $H=K_{\lfloor h/2\rfloor,\lceil h/2\rceil}$, so $q=q_h$, and put the doubles in the columns
$p_j=1+(j-1)(2W_h+2)$. The last one is in column $1+(q_h-1)(2W_h+2)=M_h$, and every column
after it is a sea. Proposition~\ref{prop:seas} gives a valid blackout with $m+q_h$ points for
every $m\ge M_h$.
\end{proof}

In the proof of the proposition we use three facts about a tilted rectangle $R$ of the
$h\times m$ grid.
\begin{enumerate}[(F1)]
\item $R$ has at most one corner in row $1$ and at most one in row $h$. Indeed, a corner in row $1$
lies in the smallest row of $R$, which holds only one corner (Lemma~\ref{lem:leftmost}), and
likewise for row $h$.
\item A corner of $R$ in row $1$ and a corner in row $h$ are opposite. If they were adjacent, the
side joining them would be $(s,h-1)$, and the other two corners would be these two plus a nonzero
vector $n$ perpendicular to it. The corner in row $1$ needs $n_2\ge0$ and the corner in row $h$
needs $n_2\le0$, so $n_2=0$. Then $n_1s=0$ forces $s=0$, a vertical side, which is impossible.
\item Two corners of $R$ in one column are opposite, since a vertical side would make $R$
axis-aligned.
\end{enumerate}

\begin{proof}[Proof of Proposition~\ref{prop:seas}]
Each double adds one point to a column, so $|Z|=m+q$.

\emph{Axis-aligned rectangles.} Only doubles have two black rows, and their pairs are distinct,
so $G_Z$ has no four-cycle. A six-cycle $r_a\,c_x\,r_b\,c_y\,r_c\,c_z$ would need doubles on
$\{a,b\}$, $\{b,c\}$, and $\{a,c\}$, a triangle in $H$. By Lemma~\ref{lem:girth}, $Z$
distinguishes the axis-aligned rectangles, and none of them is hidden, since a hidden one would
give a four-cycle.

\emph{Locality.} A tilted rectangle $R$ has span at most $W_h$ (Lemma~\ref{lem:span}), so its
columns meet at most one double. If they meet none, all black points in these columns lie in one
boundary row, and $R$ has at most one black corner by (F1). Then $R$ shows at least three
corners. Otherwise we shift the columns so that the double is column $0$, with black rows
$a<b$. Reflecting the rows if necessary, which exchanges rows $1$ and $h$, we may assume that the
seas in the columns of $R$ are black in row $1$ to the left of column $0$ and in row $h$ to the
right of it.

Let $K$ be the set of black corners of $R$. By (F1), $K$ contains at most one point
$u=(x_1,1)$ with $x_1<0$, at most one point $v=(x_2,h)$ with $x_2>0$, and points $d=(0,c)$
with $c\in\{a,b\}$. A black corner $(0,1)$ or $(0,h)$ counts as a point $d$, and then (F1)
excludes $u$ or $v$, respectively. We consider $|K|\ge2$.

\begin{enumerate}[(a)]
\item $K=\{u,v\}$. These are opposite by (F2). The other diagonal has midpoint
$\mu=\bigl(\frac{x_1+x_2}2,\frac{1+h}2\bigr)$ and endpoints $\mu\pm w$ with
$|w|^2=\frac14\bigl((x_2-x_1)^2+(h-1)^2\bigr)$. By (F1) its endpoints lie in rows $2$ to $h-1$,
so $|w_2|\le\frac{h-3}2$ and $w_1^2>\frac14(x_2-x_1)^2$. Hence one endpoint lies in a column
smaller than $x_1$ and the other in a column larger than $x_2$. Neither is in column $0$, so both
are white.
\item $K=\{(0,a),(0,b)\}$. These are opposite by (F3). The other diagonal has midpoint
$\bigl(0,\frac{a+b}2\bigr)$, length $b-a$, and $w_1\ne0$, so $|w_2|<\frac{b-a}2$. Its endpoints
lie in the columns $\pm w_1$ and strictly between rows $a$ and $b$, so both are white.
\item $K=\{u,d\}$, opposite. Here $c\ge2$ by (F1). The other diagonal has midpoint
$\bigl(\frac{x_1}2,\frac{1+c}2\bigr)$ and $|w|^2=\frac14\bigl(x_1^2+(c-1)^2\bigr)$. Its endpoints
are not in row $1$ by (F1), not in row $h$ by (F2), since the opposite corner of $u$ is $d$, and
not in column $0$ by (F3). So both are white. Since their rows are at least $2$, we get
$|w_2|\le\frac{c-3}2$. Thus $c\ge3$ and $w_1^2>\frac14x_1^2$, so one endpoint lies in a column
larger than $0$ and the other in a column smaller than $x_1$.
\item $K=\{u,d\}$, adjacent. The side $d-u=(|x_1|,c-1)$ has $c\ge2$, and the other corners are
$u+n$ and $d+n$ with $n=t\bigl(-(c-1),|x_1|\bigr)$ for some $t\ne0$. The corner $u+n$ lies in row
$1+t|x_1|$, so $t>0$. Then both corners lie in negative columns and below row $1$, so both are
white. They differ by $(|x_1|,c-1)$, so they form an oblique pair.
\item $K$ contains $v$ and a point $d$ but not $u$. This is (c) or (d) after the half turn
$(x,j)\mapsto(-x,h+1-j)$, which exchanges the two seas and fixes column $0$.
\item $|K|\ge3$. If $K$ contains $(0,a)$ and $(0,b)$, the other two corners lie strictly between
rows $a$ and $b$ by (b), so they are neither $u$ nor $v$. Otherwise $K=\{u,d,v\}$. Then $u$ and
$v$ are opposite by (F2), and $d$ is adjacent to both, so $(u-d)\cdot(v-d)=0$, that is,
$x_1x_2=(c-1)(h-c)$. The left side is negative and the right side is not, a contradiction.
\end{enumerate}
Hence no tilted rectangle is hidden, and every weak tilted rectangle shows exactly two corners.
In (a) to (c) and their reflections, the shown corners lie in columns on both sides of column $0$.
In (d) and its reflection, they form an oblique pair on one side of column $0$.

\emph{Axis-aligned rivals.} By Lemma~\ref{lem:rival}, only a horizontal or an oblique pair needs
checking. A horizontal pair $\{(x,j),(z,j)\}$ arises only in (a) to (c), with $x<0<z$. There
$B_x=\{1\}$ and $B_z=\{h\}$ are disjoint, so it has no rival. For an oblique pair
$\{(x,j),(z,k)\}$, a rival needs $(x,k)$ and $(z,j)$ black. In (a) to (c), the rows $j,k$ lie
in $2,\dots,h-1$ and $x,z\ne0$, so these points are white. In (d), $x,z<0$ and $j,k\ge2$, so
again they are white.

\emph{Two tilted rectangles.} Let weak tilted rectangles $R\ne R'$ show the same pair $P$. The
columns of each contain $P$ and span at most $W_h$, so their doubles are at most $2W_h$ apart.
Hence they share the double in column $0$. If $P$ is a diagonal of both, their black diagonals
have the same midpoint and the same length. Comparing the cases, these determine the black
diagonal. For example, (a) against (c) forces $c=h$ from the midpoint rows, then $x_1'=x_1+x_2$
and $|x_1'|=x_2-x_1$, so $x_1=0$ or $x_2=0$. Case (b) has its midpoint in column $0$, and (a) can
match it only with $x_1=-x_2$ and $(2x_2)^2+(h-1)^2=(b-a)^2\le(h-1)^2$, which is impossible. The
other comparisons are similar or immediate. If $P$ is a side of exactly one of them, it lies on
one side of column $0$ by (d) and on both sides by (a) to (c), a contradiction.

Suppose finally that $P$ is a side of both, so the black sides $Y=P-n$ and $Y'=P-n'$ satisfy
$Y'=Y+(n-n')$. If both are of type (d), with $Y=\{u,d\}$ and $Y'=\{u',d'\}$, then matching $u$
with $u'$ gives $n=n'$, since $d,d'$ lie in column $0$ and $u,u'$ in row $1$. Matching $u$ with
$d'$ would put $u'$ in column $-x_1>0$. Two sides of the reflected type are symmetric. If
$Y=\{u,d\}$ and $Y'=\{d',v'\}$, matching $d$ with $d'$ would put $v'$ in column $x_1<0$. Matching
$u$ with $d'$ gives $n_1-n'_1=|x_1|>0$, but $n_1<0<n'_1$, since the other two corners of $R$ lie
to the left of column $0$ and those of $R'$ to the right. So $R=R'$.

Presentations of three or more points are unique (Lemma~\ref{lem:threecorners}), so $Z$ is
valid.
\end{proof}

By Theorem~\ref{thm:general}, $k_r=q_h$ for every residue in Theorem~\ref{thm:periodic}(a). The
threshold $M_h$ is far from the truth for small $h$. It is $61$ for $h=5$, while
$f(5,m)=m+6$ already from $m=8$, and the least widths for $h=4$, $5$, $6$, and $7$ are $6$, $8$,
$15$, and at most $24$.
